\subsection{Discussion of Experimental Results}

This chapter provides a critical interpretation of the experimental findings presented in Chapter~5, examining the results from the perspective of the five research objectives, the identified research gaps, and the practical implications for industrial deployment. The discussion synthesizes the quantitative evidence from four experimental components --- T-GCN forecasting, supervised classification and regression, non-quantum agentic pipeline, and quantum-augmented pipeline --- and relates the findings to the broader literature on predictive maintenance.

\subsection{Spatiotemporal Forecasting: T-GCN vs.\ LSTM Baselines}

The T-GCN with cross-correlation adjacency achieves the best average R\textsuperscript{2} of 0.7702 across five sensor channels, outperforming both the OptimizedLSTM (R\textsuperscript{2} = 0.6853, $\Delta = +12.4\%$) and the Hybrid LSTM-GCN (R\textsuperscript{2} = 0.7134, $\Delta = +8.0\%$). This result confirms the core hypothesis underlying the T-GCN design: that modeling the spatial dependencies between machines through a cross-correlation lag graph provides information that temporal-only recurrent architectures cannot recover from the time series alone.

The improvement is most pronounced for humidity prediction (R\textsuperscript{2}: 0.9080 vs.\ 0.8761 for LSTM, $\Delta = +3.6\%$) and energy consumption (R\textsuperscript{2}: 0.6581 vs.\ 0.5522, $\Delta = +19.2\%$). The energy channel improvement is particularly notable because energy consumption exhibits near-zero linear correlation with other sensors individually, yet benefits substantially from graph-based cross-machine information propagation. This suggests that energy patterns across machines carry latent correlated structure --- possibly reflecting shared electrical load or production scheduling dependencies --- that the cross-correlation graph successfully captures.

The pressure sensor channel is the most difficult to forecast across all three models (R\textsuperscript{2} range: 0.4254--0.5002). Pressure has only 401 unique values in the dataset, suggesting quantized setpoint-driven dynamics that follow discrete state transitions rather than smooth continuous degradation trajectories. These step-function dynamics are inherently harder to model with gradient-based temporal architectures trained on mean-squared error, and future work should investigate specialized architectures for semi-discrete sensor channels.

The Hybrid LSTM-GCN achieves intermediate performance (R\textsuperscript{2} = 0.7134) with the shortest training time (47.89 seconds), offering a practical trade-off when computational resources are limited. However, its advantage over the pure LSTM is modest ($\Delta = +4.1\%$), suggesting that the naive concatenation of LSTM and GCN outputs without the structured lag-based graph topology of the T-GCN fails to fully realize the benefit of graph-based cross-machine reasoning. The T-GCN's principled construction of directed, lag-weighted edges from temporal cross-correlations is the key distinguishing design choice that delivers the superior performance.

Training time for the T-GCN (111.28 seconds) is approximately twice that of the LSTM (52.56 seconds), attributable to the graph convolution operations over 538 edges at each forward pass. In an operational setting where models are retrained periodically rather than continuously, this overhead is acceptable given the performance benefit. The \texttt{LearningLoop}'s retraining trigger, which requires a minimum of 5 feedback records before initiating retraining, provides a natural mechanism for controlling retraining frequency in proportion to the rate of observed distributional shift.

\subsection{Supervised Classification and Regression: Ensemble Models and TabNet}

The supervised learning experiment reveals a clear stratification of task difficulty across the six prediction targets. The anomaly flag and downtime risk binary classification targets are near-perfectly discriminable by all tested models, with LightGBM achieving F1 = 0.9990 and AUC-ROC = 0.9999. These results warrant careful interpretation: near-perfect classification performance suggests that these labels may be deterministically derivable from other features in the dataset --- specifically, the strong correlation ($r \approx 1.0$) between anomaly flag and downtime risk implies that they are effectively the same binary signal. In a real-world deployment context, a model that near-perfectly predicts anomaly flag or downtime risk from the same sensor reading may be learning a data leakage relationship rather than a genuinely predictive one. Future work should verify the temporal independence of these labels relative to the input features.

Machine status classification (best F1 = 0.712, XGBoost) and maintenance required classification (best F1 = 0.872, XGBoost) present more meaningful learning challenges that are representative of real-world difficulty. The failure type five-class classification (best F1 = 0.880, XGBoost) is the most challenging classification task, reflecting the severe class imbalance between Normal operation (91.90\%) and minority failure types (1.01--3.13\%). XGBoost's consistent superiority for the most challenging classification targets (machine status, failure type, maintenance required) is consistent with the broader literature demonstrating gradient-boosted tree advantages on tabular data with class imbalance and mixed feature types.

TabNet consistently underperforms relative to gradient-boosted ensemble methods across all targets, consistent with recent empirical comparisons that show TabNet's attention-based selection mechanism provides less benefit on tabular datasets with moderate feature counts and clear feature importances. The 50-machine dataset, with only five raw sensor channels and a modest set of engineered features, may not provide sufficient feature complexity to benefit from the sequential attention architecture that distinguishes TabNet.

The RUL regression results (best R\textsuperscript{2} = 0.189, CatBoost) deserve particular scrutiny. An R\textsuperscript{2} of 0.189 indicates that the five sensor channels and derived features in a single-timestep snapshot explain only 18.9\% of the variance in the RUL label. This is expected: the near-uniform RUL distribution (mean 234.27 hours, standard deviation 150.06 hours) means that a reading with, say, temperature = 75\textdegree C and vibration = 50 could correspond to any RUL value between 1 and 499 hours depending on the machine's degradation history and initial condition --- information that is absent from a single-snapshot feature vector. The T-GCN forecasting architecture addresses this fundamental limitation by modeling the temporal trajectory of sensor evolution, which carries the degradation history information that static snapshots lack. This result reinforces the methodological rationale for the T-GCN component in the QASAMAP pipeline. It is important to note that the T-GCN was evaluated for sensor forecasting performance, not for direct RUL prediction; extending the T-GCN architecture to predict RUL from the forecasted sensor trajectories remains a direction for future work.

\subsection{Layer 2 Agentic AI Multi-Tier Detection Empirical Evaluation}

The Layer~2 agentic AI ensemble is empirically evaluated in Chapter~5 through the pre-registered E-L2 protocol against composite ground truth (Lei 2018 + ISO 13374-2) on 25 test machines $\times$ 5 windows = 125 cases. The active V3 ensemble (3 cloud LLMs: \texttt{glm-5.1:cloud}, \texttt{kimi-k2.6:cloud}, \texttt{deepseek-v4-pro:cloud}, majority-vote with tie-break to highest tier) achieves Quadratic-Weighted Kappa $\kappa = 0.7425$ (95\% CI [0.619, 0.835]) --- substantial agreement per Landis-Koch 1977 --- with both point estimate and 95\% CI lower bound exceeding the pre-registered $\kappa \geq 0.6$ threshold required to authorise downstream Layer~3 scheduling.

The validation protocol evolved through three pre-registered methodology deviations, all transparently documented: V1 (sensor-only prompt, 4 LLMs) yielded $\kappa = 0.362$ STOP; V2 (decision-tree prompt, 3 LLMs after dropping qwen3.5 due to 53\% empty-content technical incompatibility) yielded $\kappa = 0.191$ STOP, exposing an information-asymmetry methodological flaw (the LLM was being asked to predict a label whose composite GT formula uses 4 fields the LLM was not given access to); V3 (realistic deployment input set including upstream classifier outputs --- RUL forecast, anomaly flag, downtime risk, maintenance flag --- alongside sensor readings and failure type) yielded the passing $\kappa = 0.7425$. Per-LLM ablation shows the ensemble outperforms the best individual member (deepseek $\kappa = 0.728$), confirming the value of multi-LLM aggregation conditional on member quality (V1 demonstrated that a noisy member can drag the ensemble below best individual).

Operational safety is established by perfect Critical-tier precision (1.000 --- when ensemble flags Critical, it IS Critical) and perfect Critical-or-High coverage of true Critical machines (1.000 --- no truly Critical machine routed to lower tiers). The 5/11 Critical $\to$ High downgrade is conservative-but-acceptable: those machines still receive same-day inspection rather than within-hours intervention. End-to-end ensemble p95 latency of 32.7~s is well within the 60-s Kafka producer cycle for real-time deployment.

\subsection{Layer 3 Scheduling Optimisation Empirical Evaluation}

The Layer~3 scheduling optimiser is empirically evaluated in Chapter~6 through two pre-registered experiments. E-L3A (static scheduling) sweeps 6 solvers $\times$ 3 problem sizes (Small 5 machines $\times$ 1 day; Medium 20 $\times$ 3; Large 50 $\times$ 5) $\times$ 3 employee groups $\times$ 5 random seeds = 315 controlled solver runs. All four classical heuristics (Greedy, Genetic Algorithm, Simulated Annealing, Tabu Search) converge to identical optimal makespan with 100\% feasibility and sub-second wall-clock at the largest tested scale (50 machines $\times$ 5 days $\times$ 20 employees = 45000 binary variables). The two evaluated quantum-class methods --- simulated Quantum Annealing via dwave-neal Ising mode (separate from classical SA) and quantum-inspired Simulated Bifurcation per Goto et al.\ 2019 (\emph{Sci.\ Adv.}) --- consistently fail to produce feasible solutions on the dense scheduling QUBO encoding (feasibility rate 0.0--0.2 vs 1.0 classical), with paired Wilcoxon signed-rank $p < 10^{-5}$ versus classical-best on 25 paired observations. Gate-model QAOA on Qiskit Aer statevector simulator is not applicable at meaningful problem sizes due to the qubit budget ($\sim$18 qubits practical, versus $\geq$45 variables required at the smallest meaningful instance).

E-L3B (dynamic rolling-horizon real-time feasibility) confirms that all four classical solvers sustain real-time feasibility within the 60-s Kafka producer cycle under event-driven re-optimisation across a simulated 8-hour QASAMAP working day (45 events/seed $\times$ 3 seeds = 135 re-optimisation events). The p95 wall-clock per re-opt is $<$ 1.5~s for all four classical solvers (Greedy $<$ 0.001~s, SA 0.12~s, GA 0.92~s, TS 1.36~s), with zero deadline-miss events observed. Combined with the Layer~2 ensemble e2e p95 latency of 32.7~s, the integrated Layer~2 + Layer~3 pipeline operates within the 60-s real-time budget by $\geq 26$ s margin for any classical solver choice.

The honest empirical finding at simulator-only evaluation is that QASAMAP-scale maintenance scheduling is currently best-served by classical heuristics; quantum advantage on this problem class remains unsupported by simulator-scale empirical evidence. Real D-Wave Advantage hardware was sought but the Leap free tier was not available at runtime; real-NISQ-hardware replication is reserved as future work. The Cerezo et al.\ 2025 quantum-machine-learning skepticism literature is acknowledged and addressed in Chapter~6.

\subsection{Implications for Industrial Deployment}

The experimental results collectively support three practical implications for industrial deployment of the QASAMAP framework. First, the T-GCN's superior forecasting performance over LSTM baselines, combined with the \texttt{AnchorBufferManager}'s ability to handle partial sensor instrumentation, makes it feasible to deploy T-GCN-based streaming inference in facilities where not all machines are fully instrumented. The imputation of missing sensor values from T-GCN predictions reduces the minimum sensor infrastructure required for high-quality streaming inference.

Second, the near-perfect performance of tree-based ensemble models on anomaly flag and downtime risk classification, combined with the meaningful performance on failure type and maintenance required classification, suggests that a hybrid decision architecture --- using fast ensemble model inference for initial screening and routing --- could substantially reduce the LLM agent invocation rate, improving system throughput and reducing operational cost. Events classified as anomaly flag $= 1$ with high confidence by LightGBM could be directly escalated without requiring the full multi-agent reasoning chain.

Third, the Layer~3 scheduling results show that classical heuristics are sufficient and optimal for QASAMAP-scale maintenance scheduling at simulator-only evaluation. All four classical solvers tested (Greedy, GA, SA, TS) produce feasible, optimal-makespan schedules within the 60-s Kafka real-time budget, with the simplest (Greedy) being both fastest and equally optimal. This means a deployed QASAMAP system can use the simplest classical heuristic with high confidence today. The quantum-class methods evaluated (simulated Quantum Annealing, quantum-inspired Simulated Bifurcation) do not currently provide value on this dense scheduling QUBO encoding at simulator-only evaluation. The architectural-readiness positioning for the NISQ-Advanced phase (2025--2028) is anchored on industrial precedents from adjacent scheduling problem classes (Ford Otosan + D-Wave 6$\times$ production scheduling speedup; BASF + D-Wave 7200$\times$ liquid-filling speedup), which demonstrate that hybrid quantum-classical workflows can deliver practical value with the right problem class and hardware. Real D-Wave Advantage hardware testing on QASAMAP scheduling is reserved as future work pending Leap access; gate-model QAOA at meaningful QASAMAP problem sizes requires fault-tolerant or far-larger-scale NISQ hardware than presently available.